\chapter{कार्बोनिल: नाभिकरागी योग, ऐसीटैल और रक्षण}\label{ch:b1:acetals-protection}

ग्लूकोस को जल में घोलिए, तो उसका लगभग कुछ भी पाठ्यपुस्तकों में खींचे गए खुली-शृंखला
ऐल्डिहाइड के रूप में नहीं रहता: अणु मुड़ जाता है, उसका अपना हाइड्रॉक्सिल
समूह उसके ऐल्डिहाइड से जुड़ जाता है, और एक छह-सदस्यीय वलय बंद हो जाता है। कार्बोनिल
समूह पर ऐल्कोहॉल का यही योग, अम्ल \omterm{def:b1:elementary-steps:catalyst}{उत्प्रेरक} के साथ एक बार और दोहराने पर,
\omterm{def:b1:acetals-protection:hemiacetal}{ऐसीटैल} देता है, ऐसे यौगिक जो क्षारकों और सबसे
क्रियाशील \omterm{def:b1:electronic-effects:nucleophile}{नाभिकरागियों} के प्रति स्थायी हैं। यह स्थायित्व एक उपकरण है: जो ऐल्डिहाइड
\omterm{def:b1:organomagnesium:organometallic}{ग्रीन्यार अभिकर्मक} से नष्ट हो जाता, उसे \omterm{def:b1:acetals-protection:hemiacetal}{ऐसीटैल} के रूप में छिपाकर
अभिक्रिया से पार ले जाया जा सकता है, और अंत में वापस पाया जा सकता है। यह अध्याय
\omterm{def:b1:electronic-effects:nucleophile}{इलेक्ट्रॉनरागी} के रूप में कार्बोनिल समूह, जल और ऐल्कोहॉलों के योगों,
और \omterm{def:b1:acetals-protection:protecting-group}{रक्षण} की रणनीति का अध्ययन करता है।

\begin{recall}
कार्बोनिल यौगिकों पर ग्रीन्यार योग (\cref{ch:b1:organomagnesium});
\omterm{def:b1:extent-q-and-k:quotient}{अभिक्रिया भागफल} $Q$ और $Q \neq K$ रहते हुए किसी निकाय का विकास
(\cref{ch:b1:extent-q-and-k}); \omterm{def:b1:electronic-effects:curly-arrow}{वक्र तीर}, \omterm{def:b1:electronic-effects:nucleophile}{नाभिकरागी} और
\omterm{def:b1:electronic-effects:nucleophile}{इलेक्ट्रॉनरागी} (\cref{ch:b1:electronic-effects})। विद्यालय के खंड (कक्षा
12) ने संश्लेषण के दौरान किसी समूह के \omterm{def:b1:acetals-protection:protecting-group}{रक्षण} का विचार प्रस्तुत किया था।
\end{recall}

\section{इलेक्ट्रॉनरागी के रूप में कार्बोनिल समूह}

C=O आबंध ऑक्सीजन की ओर ध्रुवित है (\cref{ch:b1:periodicity}), और
एक \omterm{def:b1:lewis-resonance-vsepr:resonance}{अनुनादी संरचना} \ce{C+-O-} कार्बन पर पूर्ण धन आवेश रखती है:
कार्बन \omterm{def:b1:electronic-effects:nucleophile}{इलेक्ट्रॉनरागी} है, और \omterm{def:b1:electronic-effects:nucleophile}{नाभिकरागी} उस पर समूह के तल के लंबवत
आक्रमण करते हैं। प्रबल \omterm{def:b1:electronic-effects:nucleophile}{नाभिकरागी} (\omterm{def:b1:organomagnesium:organometallic}{ग्रीन्यार अभिकर्मक}, हाइड्राइड,
सायनाइड) सीधे जुड़ जाते हैं। दुर्बल \omterm{def:b1:electronic-effects:nucleophile}{नाभिकरागियों} --- जल, ऐल्कोहॉल --- को सहायता चाहिए।

\begin{definition}[इलेक्ट्रॉनरागी सक्रियण]\label{def:b1:acetals-protection:activation}
किसी कार्बोनिल समूह का \emph{इलेक्ट्रॉनरागी सक्रियण}\index{इलेक्ट्रॉनरागी सक्रियण}
ऑक्सीजन पर उसका प्रोटॉनन (या किसी \omterm{def:b1:lewis-resonance-vsepr:lewis-acid}{लुइस अम्ल} से उसका उपसहसंयोजन) है।
धनायन \ce{C=OH+}, जिसका धन आवेश \omterm{def:b1:lewis-resonance-vsepr:resonance}{अनुनादी संरचना} \ce{C+-OH} के माध्यम से कार्बन के साथ
साझा होता है, दुर्बल \omterm{def:b1:electronic-effects:nucleophile}{नाभिकरागियों} के आक्रमण के लिए भी
सुलभ है।
\end{definition}

\begin{omfigure}
\setchemfig{atom sep=2.1em}
\schemestart
\chemfig{C(-[:150]R)(-[:210]H)=O}
\arrow{->[\ce{H+}]}
\chemfig{C(-[:150]R)(-[:210]H)=O^{+}-H}
\arrow{<->}
\chemfig{C^{+}(-[:150]R)(-[:210]H)-O-H}
\schemestop

{\small अम्ल द्वारा किसी ऐल्डिहाइड का सक्रियण: ऑक्सीजन का प्रोटॉनन
एक ऐसा धनायन बनाता है जिसकी \omterm{def:b1:lewis-resonance-vsepr:resonance}{अनुनादी संरचना} आवेश को
कार्बन पर रखती है।}
\end{omfigure}

\section{हाइड्रेट और हेमीऐसीटैल}

जल ऐल्डिहाइडों और कीटोनों से उत्क्रमणीय रूप से जुड़ता है, \ce{R2C=O + H2O <=>
R2C(OH)2}; कीटोनों के लिए हाइड्रेट प्रायः गौण होता है और
मेथेनल के लिए प्रमुख हो सकता है। ऐल्कोहॉल भी इसी प्रकार जुड़ता है।

\begin{definition}[हेमीऐसीटैल और ऐसीटैल]\label{def:b1:acetals-protection:hemiacetal}
\emph{हेमीऐसीटैल}\index{हेमीऐसीटैल} ऐसा यौगिक है जिसमें एक कार्बन
OH और OR दोनों समूह वहन करता है, और जो किसी कार्बोनिल समूह पर ऐल्कोहॉल के योग से
बनता है। \emph{ऐसीटैल}\index{ऐसीटैल} ऐसा यौगिक है जिसमें
एक कार्बन दो OR समूह वहन करता है, और जो किसी हेमीऐसीटैल और ऐल्कोहॉल के दूसरे
अणु से, जल की हानि के साथ, बनता है।
\end{definition}

खुली-शृंखला \omterm{def:b1:acetals-protection:hemiacetal}{हेमीऐसीटैल} ऐल्डिहाइड और ऐल्कोहॉल के साथ अपने साम्य में प्रायः
गौण होते हैं; जब OH और C=O एक ही अणु के हों और
पाँच- या छह-सदस्यीय वलय बंद हो सके, तो चक्रीय
\omterm{def:b1:acetals-protection:hemiacetal}{हेमीऐसीटैल} प्रभावी होता है। ग्लूकोस की, और अधिक सरल
5-हाइड्रॉक्सीपेंटेनल की, यही स्थिति है।

\begin{omfigure}
\setchemfig{atom sep=2em}
\schemestart
\chemfig{HO-[:30]-[:-30]-[:30]-[:-30]-[:30]C(=[:90]O)-[:-30]H}
\arrow{<=>}
\chemfig{*6(O-(-[:-30]OH)-----)}
\schemestop

\medskip

{\small 5-हाइड्रॉक्सीपेंटेनल और उसका चक्रीय \omterm{def:b1:acetals-protection:hemiacetal}{हेमीऐसीटैल}, ऑक्सीजन युक्त एक छह-सदस्यीय
वलय। कार्बन 5 का OH कार्बन 1 के ऐल्डिहाइड से जुड़ गया है;
वलय रूप प्रभावी है, जैसे ग्लूकोस के लिए।}
\end{omfigure}

\section{ऐसीटैल}

\begin{proposition}[ऐसीटैल बनने की क्रियाविधि]\label{prop:b1:acetals-protection:acetal-mechanism}
अम्ल में, कोई ऐल्डिहाइड या कीटोन और कोई ऐल्कोहॉल पहले \omterm{def:b1:acetals-protection:hemiacetal}{हेमीऐसीटैल}, फिर
\omterm{def:b1:acetals-protection:hemiacetal}{ऐसीटैल} देते हैं, इन चरणों से: C=O का प्रोटॉनन; ऐल्कोहॉल का योग;
एक प्रोटॉन की हानि (\omterm{def:b1:acetals-protection:hemiacetal}{हेमीऐसीटैल}); OH का प्रोटॉनन; जल की हानि,
जिससे शेष OR समूह द्वारा स्थायी एक \omterm{def:b1:electronic-effects:intermediates}{कार्बधनायन} (एक ऑक्सोकार्बेनियम
आयन) बनता है; दूसरे ऐल्कोहॉल का योग; एक प्रोटॉन की हानि। हर चरण
उत्क्रमणीय है।
\end{proposition}

\begin{proof}
हर चरण या तो अम्ल--क्षारक प्रोटॉन स्थानांतरण है या किसी \omterm{def:b1:electronic-effects:nucleophile}{इलेक्ट्रॉनरागी} कार्बन पर
किसी \omterm{def:b1:electronic-effects:nucleophile}{नाभिकरागी} का योग (या उससे हानि), जैसा
\cref{ch:b1:electronic-effects} में वर्णित है। जल की हानि इसलिए संभव है क्योंकि
बने धनायन, $\mathrm{R_2C^+{-}OR}$, की एक \omterm{def:b1:lewis-resonance-vsepr:resonance}{अनुनादी संरचना} $\mathrm{R_2C{=}O^+R}$ है
जिसमें हर परमाणु का अष्टक पूरा है (\cref{prop:b1:electronic-effects:carbocation-order})।
अम्ल प्रोटॉननों में खपता है और विप्रोटॉननों में पुनर्जनित होता है:
वह एक \omterm{def:b1:elementary-steps:catalyst}{उत्प्रेरक} है।
\end{proof}

\begin{omfigure}
\setchemfig{atom sep=1.9em}
\schemestart
\chemfig{C(-[:150]R)(-[:210]H)=O}
\arrow{->[\ce{H+}][]}
\chemfig{C^{+}(-[:150]R)(-[:210]H)-OH}
\arrow{->[$\mathrm{R'OH}$][\ce{-H+}]}
\chemfig{C(-[:150]R)(-[:210]H)(-[2]OR')-OH}
\schemestop

\medskip
\schemestart
\chemfig{C(-[:150]R)(-[:210]H)(-[2]OR')-OH}
\arrow{->[\ce{H+}][\ce{-H2O}]}
\chemfig{C^{+}(-[:150]R)(-[:210]H)-OR'}
\arrow{->[$\mathrm{R'OH}$][\ce{-H+}]}
\chemfig{C(-[:150]R)(-[:210]H)(-[2]OR')-OR'}
\schemestop
\par\vspace{6pt}

{\small \omterm{def:b1:acetals-protection:hemiacetal}{ऐसीटैल} का बनना, दो पंक्तियों में: पहले \omterm{def:b1:acetals-protection:hemiacetal}{हेमीऐसीटैल} (सक्रियण,
ऐल्कोहॉल का योग, एक प्रोटॉन की हानि), फिर \omterm{def:b1:acetals-protection:hemiacetal}{ऐसीटैल} (OH का प्रोटॉनन,
स्थायी धनायन तक जल की हानि, दूसरे
ऐल्कोहॉल का योग, एक प्रोटॉन की हानि)।}
\end{omfigure}

समग्र अभिक्रिया, $\mathrm{RCHO} + 2\,\mathrm{R'OH} \rightleftharpoons \mathrm{RCH(OR')_2} + \ce{H2O}$, का
साम्य स्थिरांक साधारण है, और जल एक उत्पाद है।

\begin{proposition}[ऐसीटैलन को आगे धकेलना]\label{prop:b1:acetals-protection:dean-stark}
यदि बनने वाला जल लगातार हटाया जाए, तो ऐसीटैलन तब तक चलता रहता है
जब तक कार्बोनिल यौगिक (या ऐल्कोहॉल) खप न जाए।
\end{proposition}

\begin{proof}
\omterm{def:b1:extent-q-and-k:quotient}{अभिक्रिया भागफल} $Q = [\text{ऐसीटैल}][\ce{H2O}]/([\text{ऐल्डिहाइड}]
[\mathrm{R'OH}]^2)$ तब तक $K$ से नीचे रहता है जब तक जल हटाया जाता है:
\cref{ch:b1:extent-q-and-k} के अनुसार अभिक्रिया आगे बढ़ती रहती है, $K$ का
मान चाहे जो हो। किसी डाइऑल (एथेन-1,2-डाइऑल) के साथ \omterm{def:b1:acetals-protection:hemiacetal}{ऐसीटैल} चक्रीय होता है, और
साम्य भी अधिक अनुकूल होता है, क्योंकि डाइऑल का एक अणु ऐल्कोहॉल के दो अणुओं का
स्थान लेता है।
\end{proof}

\begin{omfigure}
\begin{tikzpicture}[font=\small]
  \pic at (0,0) {heatingmantle};
  \pic at (0,0) {rbflask={0.4}{omDef!14}};
  \pic (d) at (0,3.0) {deanstark={0.35}{omDef!35}{orange!15}};
  \pic at (0,6.4) {condenser};
  \draw[->] (2.4,4.0) node[right, align=left] {संग्राहक: जल (निचली परत)\\आगे नहीं लौटता; टॉलूईन\\उफनकर फ़्लास्क में लौटता है} -- (0.8,3.8);
  \draw[->] (-2.2,1.0) node[left, align=right] {ऐल्डिहाइड, डाइऑल, टॉलूईन,\\अम्ल उत्प्रेरक, उबलता हुआ} -- (-0.7,0.9);
  \draw[->] (1.8,8.0) node[right] {संघनित्र} -- (0.4,7.8);
\end{tikzpicture}

{\small एक डीन--स्टार्क जाल। टॉलूईन और जल की वाष्प संघनित होकर
अंशांकित संग्राहक में गिरती है; जल, जो अधिक घना है और टॉलूईन के साथ \omterm{def:b1:intermolecular-forces-solvents:miscibility}{मिश्रणीय} नहीं,
नीचे बैठकर रुका रहता है, और टॉलूईन उफनकर फ़्लास्क में लौट जाता है। एकत्रित
जल का आयतन \omterm{def:b1:extent-q-and-k:extent}{अभिक्रिया की प्रगति} को मापता है।}
\end{omfigure}

\begin{proposition}[ऐसीटैलों का स्थायित्व]\label{prop:b1:acetals-protection:acetal-stability}
\omterm{def:b1:acetals-protection:hemiacetal}{ऐसीटैल} उदासीन और क्षारकीय माध्यमों में तथा \omterm{def:b1:electronic-effects:nucleophile}{नाभिकरागियों},
\omterm{def:b1:organomagnesium:organometallic}{ग्रीन्यार अभिकर्मकों} और हाइड्राइड अभिकर्मकों के प्रति स्थायी हैं; जलीय अम्ल उनका जल-अपघटन करके
उन्हें वापस कार्बोनिल यौगिक और ऐल्कोहॉल में बदल देता है।
\end{proposition}

\begin{proof}
\omterm{def:b1:acetals-protection:hemiacetal}{ऐसीटैल} कार्बन पर ऐसा कोई \omterm{def:b1:electronic-effects:nucleophile}{निर्गामी समूह} नहीं है जिसे कोई क्षारक या \omterm{def:b1:electronic-effects:nucleophile}{नाभिकरागी}
बाहर निकाल सके: \ce{RO-}, जो \omterm{def:b1:predominant-reaction:strong-acid}{प्रबल क्षारक} है, नहीं निकलता, और आक्रमण के लिए कोई C=O
शेष नहीं है। जलीय अम्ल में ऊपर वाली क्रियाविधि उलटी चलती है: जल का
बड़ा आधिक्य उत्क्रम अभिक्रिया के लिए $Q < K$ बना देता है।
\end{proof}

\section{रक्षण और विरक्षण}

\begin{definition}[रक्षी समूह]\label{def:b1:acetals-protection:protecting-group}
\emph{रक्षी समूह}\index{रक्षी समूह} वह समूह है जिसे किसी अणु में
किसी क्रियात्मकता को अस्थायी रूप से ढकने के लिए डाला जाता है, ताकि वह अणु के
किसी अन्य भाग पर लक्षित चरण में अभिक्रिया न करे। उसे डालना
\emph{रक्षण}\index{रक्षण} है, और हटाना
\emph{विरक्षण}\index{विरक्षण}।
\end{definition}

\begin{method}[रक्षण, अभिक्रिया, विरक्षण]\label{met:b1:acetals-protection:protect}
\begin{enumerate}
  \item उस क्रियात्मकता को पहचानिए जो नियोजित चरण में अभिक्रिया करेगी, या
        अभिकर्मक को नष्ट कर देगी।
  \item ऐसा \omterm{def:b1:acetals-protection:protecting-group}{रक्षी समूह} चुनिए जो उस चरण में स्थायी हो और ऐसी परिस्थितियों में हटाया जा सके
        जिन्हें अणु का शेष भाग सह सके (क्षारकों, \omterm{def:b1:electronic-effects:nucleophile}{नाभिकरागियों} या हाइड्राइडों का
        सामना करने वाले ऐल्डिहाइड या कीटोन के लिए: एक चक्रीय \omterm{def:b1:acetals-protection:hemiacetal}{ऐसीटैल})।
  \item \omterm{def:b1:acetals-protection:protecting-group}{रक्षण} कीजिए; नियोजित चरण पूरा कीजिए; \omterm{def:b1:acetals-protection:protecting-group}{विरक्षण} कीजिए।
  \item कीमत गिनिए: दो और चरण, हर एक की अपनी \omterm{def:b1:extent-q-and-k:fractional-extent}{लब्धि}।
\end{enumerate}
\end{method}

\begin{example}[बीच में आता एक कीटोन]\label{ex:b1:acetals-protection:ketoester}
किसी ऐसे अणु के ऐल्डिहाइड समूह पर \omterm{def:b1:organomagnesium:organometallic}{ग्रीन्यार अभिकर्मक} जोड़ने के लिए जो
एक कीटोन भी वहन करता है, कोई सरल \omterm{def:b1:acetals-protection:hemiacetal}{ऐसीटैल} विकल्प नहीं है (ऐल्डिहाइड कीटोनों की तुलना में
अधिक आसानी से \omterm{def:b1:acetals-protection:hemiacetal}{ऐसीटैल} बनाते हैं, इसलिए गलत समूह रक्षित होगा);
चरणों का क्रम या अभिकर्मक बदलने होंगे। \omterm{def:b1:acetals-protection:protecting-group}{रक्षण} एक
रणनीति है, कोई प्रतिवर्त क्रिया नहीं।
\end{example}

\section{अभ्यास}

\begin{exercise}[$\star$]\label{exo:b1:acetals-protection:1}
इनमें \omterm{def:b1:acetals-protection:hemiacetal}{हेमीऐसीटैल}, \omterm{def:b1:acetals-protection:hemiacetal}{ऐसीटैल}, हाइड्रेट, ईथर या ऐल्कोहॉल कार्बन पहचानिए:
\ce{CH3CH(OH)OCH3}, \ce{CH3CH(OCH3)2}, \ce{CH3CH(OH)2}, \ce{CH3OCH3},
\ce{(CH3)2C(OCH2CH3)2}।
\end{exercise}

\begin{exercise}[$\star$]\label{exo:b1:acetals-protection:2}
मेथेनॉल के साथ प्रोपेनल के \omterm{def:b1:acetals-protection:hemiacetal}{ऐसीटैल} के, और एथेन-1,2-डाइऑल के साथ
प्रोपेनोन के चक्रीय \omterm{def:b1:acetals-protection:hemiacetal}{ऐसीटैल} के बनने का समग्र समीकरण लिखिए।
\end{exercise}

\begin{exercise}[$\star$]\label{exo:b1:acetals-protection:3}
इनमें से कौन-से अभिकर्मक \omterm{def:b1:acetals-protection:hemiacetal}{ऐसीटैल} को अपरिवर्तित छोड़ते हैं: सोडियम हाइड्रॉक्साइड
विलयन, मेथिलमैग्नीशियम ब्रोमाइड, जल में तनु सल्फ़्यूरिक अम्ल, सोडियम
बोरोहाइड्राइड?
\end{exercise}

\begin{exercise}[$\star$]\label{exo:b1:acetals-protection:4}
4-हाइड्रॉक्सीब्यूटेनल द्वारा बना चक्रीय \omterm{def:b1:acetals-protection:hemiacetal}{हेमीऐसीटैल} खींचिए। वलय का आकार
क्या है?
\end{exercise}

\begin{exercise}[$\star\star$]\label{exo:b1:acetals-protection:5}
एथेनल के डाइमेथिल \omterm{def:b1:acetals-protection:hemiacetal}{ऐसीटैल} के अम्ल-उत्प्रेरित बनने की पूरी क्रियाविधि
\omterm{def:b1:electronic-effects:curly-arrow}{वक्र तीरों} के साथ लिखिए, और दिखाइए कि अम्ल
पुनर्जनित होता है।
\end{exercise}

\begin{exercise}[$\star\star$]\label{exo:b1:acetals-protection:6}
2,2-डाइमेथॉक्सीप्रोपेन के अम्लीय जल-अपघटन की क्रियाविधि लिखिए।
जल की बड़ी मात्रा क्यों प्रयुक्त की जाती है?
\end{exercise}

\begin{exercise}[$\star\star$]\label{exo:b1:acetals-protection:7}
किसी कीटोन के \qty{0.250}{mol} के ऐसीटैलन का अनुसरण
डीन--स्टार्क जाल से किया जाता है। पूर्ण रूपांतरण पर जल का कितना आयतन एकत्रित होना चाहिए
(घनत्व लगभग \qty{1.0}{g/mL})? एक घंटे बाद
\qty{3.2}{mL} एकत्रित हुआ है: रूपांतरण कितना है?
\end{exercise}

\begin{exercise}[$\star\star$]\label{exo:b1:acetals-protection:8}
समझाइए कि एक ही कार्बोनिल यौगिक के लिए एथेन-1,2-डाइऑल के साथ चक्रीय \omterm{def:b1:acetals-protection:hemiacetal}{ऐसीटैल}
मेथेनॉल के दो अणुओं वाले \omterm{def:b1:acetals-protection:hemiacetal}{ऐसीटैल} की तुलना में अधिक पूर्णता से क्यों
बनता है।
\end{exercise}

\begin{exercise}[$\star\star$]\label{exo:b1:acetals-protection:9}
किसी रक्षित यौगिक को \omterm{def:b1:organomagnesium:organometallic}{ग्रीन्यार अभिकर्मक} से उपचारित करने से पहले अम्ल \omterm{def:b1:elementary-steps:catalyst}{उत्प्रेरक} को
हटाना (या उदासीन करना) क्यों आवश्यक है?
\end{exercise}

\begin{exercise}[$\star\star\star$]\label{exo:b1:acetals-protection:10}
एक \omterm{def:b1:acetals-protection:protecting-group}{रक्षण} का उपयोग करते हुए, 4-क्लोरोब्यूटेन-2-ओन और प्रोपेनोन से
5-हाइड्रॉक्सी-5-मेथिलहेक्सेन-2-ओन, \ce{(CH3)2C(OH)CH2CH2COCH3}, का एक संश्लेषण
प्रस्तावित कीजिए। हर चरण, उसके अभिकर्मक और हर एक की भूमिका बताइए।
\end{exercise}

\begin{exercise}[$\star\star\star$]\label{exo:b1:acetals-protection:11}
\omterm{def:b1:extent-q-and-k:quotient}{अभिक्रिया भागफल} से दिखाइए कि मेथेनॉल का बड़ा आधिक्य भी
किसी ऐल्डिहाइड के ऐसीटैलन को आगे धकेलता है, और इस विधि की तुलना
जल हटाने से कीजिए।
\end{exercise}

\begin{exercise}[$\star\star\star$]\label{exo:b1:acetals-protection:12}
समझाइए कि अम्लीय मेथेनॉल में ग्लूकोस के चक्रीय \omterm{def:b1:acetals-protection:hemiacetal}{हेमीऐसीटैल} पर आक्रमण होकर
एक मेथिल \omterm{def:b1:acetals-protection:hemiacetal}{ऐसीटैल} (एक मेथिल ग्लाइकोसाइड) क्यों बनता है, जबकि ग्लूकोस के अन्य OH
समूह ईथरों में नहीं बदलते।
\end{exercise}

\section{समस्या: ऐल्डिहाइड वहन करने वाला एक ग्रीन्यार अभिकर्मक}

\begin{problem}[{सप्ताहांत समस्या --- 4-ब्रोमोब्यूटेनल ग्रीन्यार अभिकर्मक क्यों नहीं
दे सकता, डीन--स्टार्क जाल से चक्रीय ऐसीटैल के रूप में उसका रक्षण,
प्रोपेनोन पर ग्रीन्यार योग, विरक्षण, और
एकत्रित जल का आयतन}]
\label{pb:b1:acetals-protection:1}
एक रसायनज्ञ को 5-हाइड्रॉक्सी-5-मेथिलहेक्सेनल, \ce{(CH3)2C(OH)CH2CH2CH2CHO}, चाहिए, और
वह इसे 4-ब्रोमोब्यूटेनल, \ce{BrCH2CH2CH2CHO}, और प्रोपेनोन से बनाने की योजना बनाती है।
वह \qty{15.09}{g} 4-ब्रोमोब्यूटेनल से आरंभ करती है। मोलर द्रव्यमान
(\unit{g/mol}): \ce{C4H7BrO} 150.9, \ce{C2H6O2} 62.0, \ce{C6H11BrO2} 194.9,
\ce{C3H6O} 58.0, \ce{C9H18O3} 174.0, \ce{C7H14O2} 130.0, \ce{H2O} 18.0;
जल का घनत्व \qty{0.997}{g/mL}।
% ledger: rho:water

\medskip
\noindent\textbf{भाग I --- समस्या।}
\begin{enumerate}
  \item वह \omterm{def:b1:organomagnesium:organometallic}{ग्रीन्यार अभिकर्मक} लिखिए जो 4-ब्रोमोब्यूटेनल देता।
  \item वह अभिकर्मक अस्तित्व में क्यों नहीं रह सकता? उसे नष्ट करने वाली अभिक्रिया
        लिखिए।
  \item किस क्रियात्मकता का \omterm{def:b1:acetals-protection:protecting-group}{रक्षण} आवश्यक है, और किसके विरुद्ध?
  \item \omterm{def:b1:acetals-protection:hemiacetal}{ऐसीटैल} उपयुक्त क्यों है?
  \item चक्रीय \omterm{def:b1:acetals-protection:hemiacetal}{ऐसीटैल} (एथेन-1,2-डाइऑल के साथ) को वरीयता क्यों दी जाती है?
  \item 4-ब्रोमोब्यूटेनल की मात्रा परिकलित कीजिए।
  \item \qty{20}{\percent} आधिक्य के लिए डाइऑल का द्रव्यमान परिकलित कीजिए।
\end{enumerate}

\medskip
\noindent\textbf{भाग II --- \omterm{def:b1:acetals-protection:protecting-group}{रक्षण}।}
\begin{enumerate}[resume]
  \item एथेन-1,2-डाइऑल के साथ ऐसीटैलन का समीकरण लिखिए।
  \item अम्ल \omterm{def:b1:elementary-steps:catalyst}{उत्प्रेरक} (4-मेथिलबेंज़ीनसल्फ़ोनिक
        अम्ल) की क्या भूमिका है?
  \item डीन--स्टार्क जाल की कार्यविधि का वर्णन कीजिए।
  \item $Q$ और $K$ से समझाइए कि जल हटाना अभिक्रिया को आगे क्यों धकेलता है।
  \item फ़्लास्क में टॉलूईन क्यों लौटाया जाता है, जल क्यों नहीं?
  \item रसायनज्ञ को कैसे पता चलता है कि अभिक्रिया पूर्ण हो गई है?
  \item पूर्ण रूपांतरण पर अपेक्षित रक्षित ब्रोमाइड का द्रव्यमान परिकलित कीजिए।
\end{enumerate}

\medskip
\noindent\textbf{भाग III --- ग्रीन्यार चरण।}
\begin{enumerate}[resume]
  \item रक्षित ब्रोमाइड से \omterm{def:b1:organomagnesium:organometallic}{ग्रीन्यार अभिकर्मक} का बनना
        लिखिए।
  \item प्रोपेनोन पर उसका योग और बनने वाला \omterm{def:b1:alcohol-activation:alkoxide}{ऐल्कॉक्साइड} लिखिए।
  \item इस चरण का उपचार अम्लीय के बजाय हल्का क्षारकीय या उदासीन (उदाहरण के लिए
        जलीय अमोनियम क्लोराइड) क्यों होना चाहिए?
  \item इस चरण की \qty{70}{\percent} \omterm{def:b1:extent-q-and-k:fractional-extent}{लब्धि} के लिए रक्षित ऐल्कोहॉल का
        द्रव्यमान परिकलित कीजिए।
\end{enumerate}

\medskip
\noindent\textbf{भाग IV --- \omterm{def:b1:acetals-protection:protecting-group}{विरक्षण}।}
\begin{enumerate}[resume]
  \item जलीय अम्ल में \omterm{def:b1:acetals-protection:protecting-group}{विरक्षण} लिखिए।
  \item तृतीयक ऐल्कोहॉल मृदु जलीय अम्ल में क्यों बचा रहता है?
  \item उत्पाद अपने OH को ऐल्डिहाइड से जोड़कर एक छह-सदस्यीय वलय बंद कर
        सकता है। वह चक्रीय \omterm{def:b1:acetals-protection:hemiacetal}{हेमीऐसीटैल} खींचिए।
  \item \omterm{def:b1:acetals-protection:protecting-group}{विरक्षण} की \qty{90}{\percent} \omterm{def:b1:extent-q-and-k:fractional-extent}{लब्धि} पर हाइड्रॉक्सी ऐल्डिहाइड का
        द्रव्यमान परिकलित कीजिए।
  \item रसायनज्ञ 4-ब्रोमोब्यूटेनल से कितनी समग्र \omterm{def:b1:extent-q-and-k:fractional-extent}{लब्धि} तक पहुँचती है?
  \item पूर्ण रूपांतरण पर \omterm{def:b1:acetals-protection:protecting-group}{रक्षण} चरण में बने जल का द्रव्यमान
        परिकलित कीजिए।
  \item पूर्ण रूपांतरण पर डीन--स्टार्क जाल में एकत्रित जल का आयतन
        बताइए।
\end{enumerate}
\end{problem}
